\documentclass[10pt]{article}
\usepackage[utf8]{inputenc}
\usepackage[T1]{fontenc}

\title{Mapping of the Logical Structure of Partial Differential Equations using an Agentic Knowledge Graph Framework }

\author{Name: Varuneesh Johar\\
Enrolment No.: 21323040\\
Supervisor: Prof. Ram Jiwari}
\date{}


\begin{document}
\maketitle


\section*{Motivation}
Partial Differential Equations (PDEs) are foundational to modern science and engineering. However, the literature in this field is vast and deeply interconnected, making it a significant challenge to discern the complex relationships between equations, theorems, and numerical methods. This reliance on manual exploration and expert intuition slows down progress and raises the barrier of entry for newcomers. To address this knowledge navigation problem, this thesis proposes a computational approach to systematically map the logical structure of PDE research.

\section*{Proposed Framework}
This thesis introduces a novel agentic Knowledge Graph (KG) framework tailored to the unique demands of mathematical literature. The approach builds upon recent advances in multi-agent AI for scientific discovery, employing a pipeline of specialised, Large Language Model (LLM)-powered agents to parse a curated corpus of PDE research papers. The agents' core tasks are to:

\begin{itemize}
  \item Extract key entities (e.g., named equations, theorems, solution techniques).
  \item Identify the logical and semantic relationships between them.
  \item Organise this information into a structured, machine-readable Knowledge Graph.
\end{itemize}

The research will concentrate on the domain of hyperbolic wave problems and their associated advanced numerical techniques, including Finite Element Methods (FEM), Discontinuous Galerkin (DG), Hybridised DG (HDG), and High-Order Hybrid (HHO) methods.

\section*{Anticipated Contributions}
The central contribution will be the design and implementation of the agentic system and the resulting Knowledge Graph that captures the logical structure of hyperbolic PDE methods. More broadly, the thesis aims to:

\begin{itemize}
  \item Deliver an automated framework for extracting and structuring mathematical knowledge.
  \item Produce a domain-specific KG interlinking PDE equations, theorems, and computational methods.
  \item Demonstrate how structured knowledge supports navigation and reasoning in advanced mathematical domains.
  \item Optionally, if time permits, develop a prototype frontend tool to showcase the practical utility of the KG.
\end{itemize}

\section*{References}
\begin{itemize}
  \item Bo, L., et al. (2024). SciAgents: A Multi-Agent System with a Knowledge Graph for Scientific Discovery. arXiv preprint.
  \item Hogan, A., et al. (2021). Knowledge graphs. ACM Computing Surveys (CSUR), 54(4), 1-37.
\end{itemize}

\end{document}